% v0.30 candidate component. Independent human reconstruction.
% Source pin: 7cbf7a0c32fca1badf05931879a37c3f8a2b03b4.
\section{An exact weaker analytic input on the selected route}
\label{sec:analytic-envelope-audit}

The pointwise window bound in AN1 is sufficient, but is stronger than the
actual analytic interface in the pinned source.  Its
\texttt{WindowHypU} requires a nonnegative envelope with a small weighted
integral.  The following proof reconstructs that interface independently.
The arithmetic construction of such envelopes remains an obligation.

\subsection{An exact positive-part formulation}

Fix one primitive target $\psi$ and its fixed positive character family,
and write $H=H_\psi$ as in AN1 and $g(u)=H(e^u)$.
For a real number $v$, put $v_+=\max(v,0)$.  Fix
$0<\theta<1$ and a fixed $\delta>0$.  The integral window condition is
\[
 \text{for every $\eta>0$ and $Y_0$, there exists
 $Y\ge\max(2,Y_0)$ such that}\qquad
 \int_Y^{Y^{1+\delta}} H(x)_+\,x^{-\theta-1}\,dx\le Y^\eta.
 \tag{AQ.1}
\]
In logarithmic coordinates it is
\[
 \forall\eta>0\ \forall U_0\ \exists U\ge\max(1,U_0),\qquad
 \int_U^{(1+\delta)U}g(u)_+ e^{-\theta u}\,du\le e^{\eta U}.
 \tag{AQ.2}
\]
The harmless lower thresholds $2$ and $1$ are interchangeable because
the starting threshold is universally quantified.

\begin{proposition}[AQ1: the literal envelope interface]
\label{claim:AQ1}
For this $g$, condition (AQ.2) is equivalent to the source's
\texttt{WindowHypU}$(g,\theta,\delta)$.
\end{proposition}
\begin{proof}
The source asks for a function $R$ on each selected interval satisfying
$R\ge0$, $g\le R$, integrability of $R(u)e^{-\theta u}$, and the
integral bound in (AQ.2) with $R$ in place of $g_+$.  Such an $R$
dominates $g_+$.  Conversely choose $R=g_+$.
The finite ramp sum defining $H$ is continuous, and its crude bound
$|H(x)|\le Mx^2$ for $x\ge1$ makes $g_+$ integrable on every bounded
interval.  This proves both implications.  Substitution $x=e^u$
proves equivalence with (AQ.1).  Notice that a bound on the signed
integral of $H$ would not imply (AQ.1).
\end{proof}

\subsection{Localization on every fixed multiplicative stretch}

\begin{lemma}[AQ2: a quantitative gamma localization]
\label{claim:AQ2}
Let $\delta>0$, put
\[
 d=\min(\delta,1),\quad m=1+\frac d2,\quad
 h=\frac{8m}{d},\quad \sigma=(h+2)^2+3,\quad \kappa=m\sigma.
\]
For $U\ge1$ let $k=\lceil\kappa U\rceil$.  There is a constant
$C_\delta$ independent of $U$ such that
\[
 \int_{(0,\infty)\setminus[U,(1+d)U]}
       \frac{u^k}{k!}e^{-(\sigma-2)u}\,du
       \le C_\delta\sigma^{-k}e^{-mU}.
 \tag{AQ.3}
\]
\end{lemma}
\begin{proof}
On the lower tail insert $e^{h(U-u)}\ge1$ and on the upper tail
insert $e^{h(u-(1+d)U)}\ge1$.  Extending both integrals gives
\[
 \frac{e^{hU}}{(\sigma+h-2)^{k+1}}
 +\frac{e^{-h(1+d)U}}{(\sigma-h-2)^{k+1}}.       \tag{AQ.4}
\]
The denominators are positive.  Multiplication by $\sigma^k$ makes
the first summand at most a fixed constant times
\[
 \exp\{[h-\kappa\log(1+(h-2)/\sigma)]U\},
\]
using $k\ge\kappa U$.  The second is at most a fixed constant times
\[
 \exp\{[-h(1+d)-\kappa\log(1-(h+2)/\sigma)]U\},
\]
using $k\le\kappa U+1$.  For $v\ge0$,
$\log(1+v)\ge v-v^2/2$; for $0\le v\le1/2$,
$-\log(1-v)\le v+v^2$.  Our choice of $\sigma$ allows the latter
inequality at $v=(h+2)/\sigma$.  Since
$(m-1)h=(1+d-m)h=4m$, the two exponents are bounded respectively by
\[
 -2m+\frac{m(h-2)^2}{2\sigma}\le-m,
 \qquad
 -2m+\frac{m(h+2)^2}{\sigma}\le-m.
\]
This proves (AQ.3), including the effect of the ceiling in $k$.
\end{proof}

\begin{theorem}[AQ3: positive integral windows exclude right-hand zeros]
\label{claim:AQ3}
If (AQ.1) holds for a fixed positive character family and fixed
$0<\theta<1$, $\delta>0$, then every constituent Dirichlet
$L$-function is nonzero on $\operatorname{Re}s>\theta$, apart from
the principal pole at $1$.
\end{theorem}
\begin{proof}
The Mellin computations AN2--AN3 give the meromorphic continuation
of the Laplace transform
\[
 \Phi(s)=\int_0^\infty g(u)e^{-su}\,du
 \qquad(\operatorname{Re}s>2)
\]
to $\operatorname{Re}s>0$.  Its poles there are exactly the zeros of
the constituent functions, after combining coincident poles, and a
zero $\rho$ has residue
\[
 \frac{b_\rho}{\rho(1-\rho)},\qquad b_\rho>0.       \tag{AQ.5}
\]
The apparent pole at $1$ has been canceled by the centering; the
nonprincipal values at $1$ are nonzero by QW14.  These poles are
locally finite, lie to the left of $\operatorname{Re}s=\sigma$,
and the residue at every real pole in $(0,1)$ is positive.  The
coefficient $b_\rho$ is a positive weighted sum of multiplicities
for the characters and their conjugates; hence there is no
cancellation even when different constituent functions vanish at
the same point.  The elementary bound $|g(u)|\le M e^{2u}$ is
uniform for $u\ge0$, with $M$ depending on the fixed family.

Use $\sigma,\kappa$ from AQ2 and put
\[
 A_k(z)=\frac{(-1)^k}{k!}\Phi^{(k)}(z)
       =\int_0^\infty g(u)\frac{u^k}{k!}e^{-zu}\,du.
\]
For every real $t$, reality of $g$ gives
\[
 |A_k(\sigma+it)|+A_k(\sigma)
 \le 2\int_0^\infty g(u)_+\frac{u^k}{k!}e^{-\sigma u}\,du.
 \tag{AQ.6}
\]
Restrict the good interval in (AQ.2) to $[U,(1+d)U]$; this only
decreases its nonnegative integral.  On this interval
\[
 \frac{u^k}{k!}e^{-(\sigma-\theta)u}
          \le(\sigma-\theta)^{-k},
\]
by the exponential series.  Outside it apply AQ2 and the crude
growth bound.  Therefore the selected $U$ and
$k=\lceil\kappa U\rceil$ satisfy
\[
 |A_k(\sigma+it)|+A_k(\sigma)
 \le2e^{\eta U}(\sigma-\theta)^{-k}+O(\sigma^{-k}).
 \tag{AQ.7}
\]
The implied constant is independent of the selected window and of
$t$.  Since $U\le k/\kappa$, the first term is at most
$2B^{-k}$ where $B=(\sigma-\theta)e^{-\eta/\kappa}$.

Suppose some pole lies to the right of $\theta$.  Set
\[
 D=\min\{\sigma-\theta,\ \inf_\rho|\sigma-\rho|\}.
\]
First suppose that some pole has horizontal distance from
$\operatorname{Re}s=\sigma$ strictly less than $D$.  Choose a real
$t$ for which $z=\sigma+it$ has a unique nearest pole $\rho$ at
distance $a<D$.  This is possible near the imaginary part of that
pole: only finitely many poles can compete inside the relevant
compact disks, and each pair with different imaginary parts
forbids only one $t$.  Distinct poles with the same imaginary part
have different distances because they lie to the left of $\sigma$.
Choose $a<r<D$ so that the closed $r$-disk about $z$ contains no
other pole.  The $r$-disk about $\sigma$ contains none.  Cauchy's
estimate, after subtracting the principal part at $\rho$, gives
\[
 A_k(z)=\frac{c}{(z-\rho)^{k+1}}+O(r^{-k}),\qquad
 A_k(\sigma)=O(r^{-k}),\quad c\ne0.
\]
Choose $\eta>0$ sufficiently small that
$B=(\sigma-\theta)e^{-\eta/\kappa}>r$.
This choice precedes the selection of arbitrarily late windows in
(AQ.2).  Inequality (AQ.7) now implies
\[
 |c|a^{-k-1}\le O(r^{-k})+2B^{-k},
\]
which is impossible as $U$, hence $k$, tends to infinity.

In the remaining case all poles have horizontal distance at least
$D$.  The existence of a pole to the right of $\theta$ forces
$D<\sigma-\theta$.  A nearest pole to $\sigma$ exists by local
finiteness in a compact disk contained in
$\operatorname{Re}s>\theta$.  It must be real, since its horizontal
distance is at least $D$ and its total distance equals $D$.
It is the unique pole at that distance, located at $\beta=\sigma-D$,
and its residue $c$ is positive.  Choose
$D<r<\sigma-\theta$ with no other pole in the closed $r$-disk.
Then
\[
 A_k(\sigma)=cD^{-k-1}+O(r^{-k}).
\]
Choose $\eta>0$ so that $B>r$ and use (AQ.7) at $t=0$.
For large $k$ its left side is $2A_k(\sigma)>0$, giving the same
impossible comparison of exponential rates.  Thus no such pole
exists in either case.  The Mellin residue identification (AQ.5)
proves the asserted nonvanishing.
\end{proof}

\begin{corollary}[AQ4: a minimal sufficient terminal contract]
\label{claim:AQ4}
For every fixed primitive $\psi$ and every
$\theta\in(1/2,1)$, suppose there is a $\delta_{\psi,\theta}>0$,
fixed before $\eta$ and the window, such that (AQ.1) holds.
It suffices to impose this for a sequence of such $\theta$
decreasing to $1/2$.  Then Dirichlet GRH follows.  The ordinary RH
case requires only $\psi=1$.
\end{corollary}
\begin{proof}
AQ3 excludes any proposed zero with real part $\beta>1/2$ by
choosing a member of the sequence below $\beta$.
The primitive functional equation reflects zeros in the open
strip on the left to right-hand zeros of the conjugate character.
For an imprimitive character the additional Euler factors
$1-\chi^*(p)p^{-s}$ are nonzero for $\operatorname{Re}s>0$.
These are the already written GH6 transitions.  This conclusion
concerns Dirichlet $L$-functions; it does not extend the scope to
arbitrary automorphic $L$-functions.
\end{proof}

\paragraph{What the contract saves, and what it does not supply.}
A pointwise bound $H(x)\le C_\theta x^\theta$ gives an integral
at most $C_\theta\delta\log Y$, which is at most $Y^\eta$ on
sufficiently late windows for every $\eta>0$.  Thus the previous
pointwise route is retained.  Conversely (AQ.1) permits short
positive spikes and does not imply a pointwise bound.  It is a
potentially useful alternate arithmetic target, not a new supplier
of the host, common family, theta data, arithmetic gap or operators.
The character family and $\delta$ are fixed before $\eta$ and the
window; constants or weights chosen anew with each $Y$ would not
satisfy the theorem as stated.

\subsection{A small discrete simplification}

\begin{lemma}[AQ5: interpolation of the original centered ramps]
\label{claim:AQ5}
For $n\ge1$ and $x\in[n,n+1]$, let $L_n(x)$ be the affine
interpolant of $H(n)$ and $H(n+1)$.  For the original target with
principal coefficient $1$,
\[
 0\le H(x)-L_n(x)
       \le\frac{(x-n)(n+1-x)}{2n}\le\frac1{8n}.
 \tag{AQ.8}
\]
Consequently for $0<\theta<1$ and $C\ge0$, bounds
$H(n)\le Cn^\theta$, $H(n+1)\le C(n+1)^\theta$ imply
$H(x)\le Cx^\theta+1/(8n)$ on that interval.
\end{lemma}
\begin{proof}
Every nonprincipal centered ramp is affine between consecutive
integers, as is the finite von Mangoldt part of the principal
ramp.  The only nonaffine term is $-x\log x$; therefore
$H''(x)=-1/x$ on $(n,n+1)$.  The difference $H-L_n$ is concave
with zero boundary values, giving its lower bound.  Subtract
$(x-n)(n+1-x)/(2n)$.  The resulting function has nonnegative
second derivative and zero boundary values, so it is at most
zero.  This proves (AQ.8).  Finally $x^\theta$ is concave, so its
affine interpolant is at most its value at $x$.
\end{proof}

In particular, if for a fixed $\lambda>1$ and each
$\theta\in(1/2,1)$ there are a fixed $C_\theta\ge0$ and arbitrarily
late windows with $H(n)\le C_\theta n^\theta$ for every integer
$n\in[Y,Y^\lambda]$, the same GRH conclusion follows.  Choose
$1<\mu<\lambda$ and $Z=\lceil Y\rceil$.  For sufficiently large
$Y$, $Z^\mu+1\le Y^\lambda$, so for every $x\in[Z,Z^\mu]$,
the two integers $\lfloor x\rfloor$ and $\lfloor x\rfloor+1$
lie in the supplied window.
AQ5 gives $H(x)\le(C_\theta+1)x^\theta$ throughout that smaller
window.  AQ3 applies with $\delta=\mu-1$.

\paragraph{Pinned-source scope.}
The source defines the original \texttt{WindowHyp} and its logarithmic
version \texttt{WindowHypU}.  Its \texttt{landau\_window} proves the
logarithmic pole criterion and \texttt{lemmaL} its character-family
zero-free consequence;
the source's power-window theorem derives the
pointwise special case for any fixed $\delta>0$.
The present gamma estimate is an independent human proof of the
same analytic implication, not a claim about a new arithmetic
window construction or a fresh Lean replay.
